\answer{ex:08:1}{(a)~$\approx1.7\cdot10^6$. (b)~$g=\frac{0.9}{\sqrt{0.19}}\,\sigma_{\text{après}}/\sigma_{\text{avant}}\approx16.5$: la réponse ne dépend que du rapport des bruits.}
\answer{ex:08:2}{(a)~Valeurs propres $-(1\pm g\beta)/\tau$; la différence s'amortit en $\tau/(1-g\beta)$: $40\ms$ pour $g\beta=0.5$ (différence des entrées amplifiée trois fois), $200\ms$ pour $0.9$ ($19$ fois). (b)~$0.905$ et $1.105$; entre les deux frontières, seule l'entrée forte gagne.}
\answer{ex:08:3}{$P(k)=e^{gu_k/\sigma}/\sum_le^{gu_l/\sigma}$, c'est-à-dire~\eqref{eq:softmax}; $g=10\ln9\approx22$.}
\answer{ex:08:4}{$\bar\Delta=1/3$, $g^*\approx3\ln(1/h)$: $21$, $14$, $9$ contre $16$--$23$, $11$ et $8$ dans le modèle.}
\answer{ex:08:5}{$g^*$ va de $16$--$23$ pour $h=0.001$ à $6$--$8$ pour $h=0.2$; l'estimation $3\ln(1/h)$ est juste à un facteur $1.5$ près pour $h\le0.05$; pour $h\gtrsim\eta/10$, la règle $Q\leftarrow Q+\eta(R-Q)$ n'a pas le temps d'assimiler ce que l'exploration a trouvé, et $g^*$ reste bloqué vers $8$.}
\answer{ex:08:6}{$T_p=\tau\ln9=44\ms$ pour $g_{\text{lo}}=0$ (modèle: $44\ms$) et $70\ms$ pour $g_{\text{lo}}=0.25$: une bouffée de $2$--$3\tau$ avec un gain proche de zéro.}
\answer{ex:08:7}{(a)~Meilleur $g$ constant: $0.925$ ($g=8$), oracle: $0.929$ ($16/8$); il n'y a presque rien à gagner. (b)~Par exemple, des époques calmes à récompenses rares ($p\in[0,0.3]$) et des époques changeantes avec $h=0.1$: $0.850$ contre $0.872$; l'ajustement en récupère environ un quart pour $\eta=0.2$ et presque tout pour $\eta=0.05$.}
